\MakeAbstract{
    机械臂已在智能装配、搬运等制造业场景中得到了广泛应用。在人工智能国家战略深入推进的背景下，赋予机械臂多任务自主执行能力与自然语言交互能力已成为智能装备领域的重要发展方向，具有降低人机协作门槛、拓展应用场景等现实意义。考虑到上述需求，本研究针对当前机械臂控制中多任务泛化能力不足以及缺乏自然语言到底层控制指令的跨模态映射等问题，利用大语言模型的泛化和自然语言交互能力，提出了从强化学习专家策略训练到轻量大语言模型微调的技术路线。主要工作内容如下：

    （1）针对机械臂多任务控制中专家策略数据的来源问题，提出了基于TQC算法与HER机制的机械臂单任务专家策略训练方案。该方案在MuJoCo仿真环境中，针对Fetch机械臂的四项基础操作任务训练独立的单任务专家策略。并通过消融实验对比仅TQC与TQC+HER的训练表现，分析HER在不同任务下的作用。实验结果表明，通过该方法训练的单任务专家策略在各自任务上均具备良好性能，HER在稀疏奖励环境中能有效加速训练收敛，但在部分任务环境中会降低策略网络的最终性能。

    （2）针对基于自然语言的专家轨迹数据集的构建问题，设计了包含轨迹采集、语义转化和复现验证清洗的数据集构建流程。该流程先基于已训练的单任务专家策略在仿真环境中执行的原始轨迹，通过三角色对话模板封装与数据语义解释，转化为可供语言模型处理的自然语言轨迹数据集；进而在原始仿真引擎中对数据集中自然语言轨迹进行转译与回放，清洗不可复现轨迹。实验结果表明，经过上述流程构建的专家轨迹数据集在大多数任务下具备良好的轨迹复现率，且清洗后的轨迹均能在仿真环境中进行复现。

    （3）针对大语言模型难以精确输出底层控制策略及推理速度慢的问题，提出了多任务混合微调、部署与仿真测试方法。该方法首先基于已构建的测试集与多个常用的基础轻量大语言模型开展多任务混合LoRA微调；再通过vLLM引擎实现推理加速和本地部署；最后构建大语言模型在仿真环境中的推理测试流程，分析多任务微调后各模型性能。实验结果表明，多任务微调后的多数语言模型已初步具备完成所有四项任务的泛化能力并具备实时控制能力，但对不同任务的适应能力存在明显差异。

}{机械臂控制，大语言模型，参数高效微调，强化学习，行为克隆}

\MakeAbstractEng{
Robotic arms have seen wide use in manufacturing, particularly in assembly and material handling. Under China's national AI development strategy, enabling these systems with multi-task autonomy and natural language command interfaces has become a practical research direction, with the potential to reduce operational expertise requirements and broaden the range of deployable scenarios. Two issues remain in current practice, however: Firstly, most control methods do not generalize well across tasks, and no established mechanism exists for bridging natural language to low-level control commands. This study used the generalization and language processing strengths of large language models to tackle these issues, building a pipeline from reinforcement learning expert policy training to lightweight LLM fine-tuning. The main work is as follows:

(1) To obtain expert policy data for multi-task control, a single-task expert policy training scheme based on TQC combined with HER was proposed. In MuJoCo, four independent policies were trained for the Reach, Push, Slide, and PickAndPlace tasks on a Fetch robotic arm. TQC alone and TQC+HER were then compared through ablation experiments. All four policies performed well on their respective tasks. HER sped up convergence under sparse rewards, but reduced final policy performance in some environments.

(2) A dataset construction pipeline was designed to turn expert trajectories into natural language. Raw trajectories from the trained policies were wrapped into three-role dialogue templates and annotated with semantic information, yielding a text-based trajectory dataset for language model training. Non-reproducible trajectories were then identified and removed via closed-loop replay in the original physics engine. The resulting dataset retained high replay fidelity across most tasks, and all remaining trajectories passed replay verification.

(3) A multi-task fine-tuning, deployment, and evaluation workflow was developed to tackle two LLM limitations: difficulty generating precise control commands and slow inference. Lightweight base models were fine-tuned with LoRA on the constructed dataset, deployed locally through vLLM for faster inference, and then evaluated in a simulated control loop. Most of the fine-tuned models picked up basic multi-task generalization and real-time control capability, though performance differed noticeably across tasks.

}{Robotic arm control, Large language models, Parameter-efficient fine-tuning, Reinforcement learning, Behavior cloning}
